% The Slide Definitions
%document
\documentclass[10pt]{beamer}
%theme
\usetheme{metropolis}
% packages
\usepackage{color}
\usepackage{listings}
\usepackage[ngerman]{babel}
\usepackage[utf8]{inputenc}
\usepackage{multicol}


% color definitions
\definecolor{mygreen}{rgb}{0,0.6,0}
\definecolor{mygray}{rgb}{0.5,0.5,0.5}
\definecolor{mymauve}{rgb}{0.58,0,0.82}

\lstset{
    backgroundcolor=\color{white},
    % choose the background color;
    % you must add \usepackage{color} or \usepackage{xcolor}
    basicstyle=\footnotesize\ttfamily,
    % the size of the fonts that are used for the code
    breakatwhitespace=false,
    % sets if automatic breaks should only happen at whitespace
    breaklines=true,                 % sets automatic line breaking
    captionpos=b,                    % sets the caption-position to bottom
    commentstyle=\color{mygreen},    % comment style
    % deletekeywords={...},
    % if you want to delete keywords from the given language
    extendedchars=true,
    % lets you use non-ASCII characters;
    % for 8-bits encodings only, does not work with UTF-8
    frame=single,                    % adds a frame around the code
    keepspaces=true,
    % keeps spaces in text,
    % useful for keeping indentation of code
    % (possibly needs columns=flexible)
    keywordstyle=\color{blue},       % keyword style
    % morekeywords={*,...},
    % if you want to add more keywords to the set
    numbers=left,
    % where to put the line-numbers; possible values are (none, left, right)
    numbersep=5pt,
    % how far the line-numbers are from the code
    numberstyle=\tiny\color{mygray},
    % the style that is used for the line-numbers
    rulecolor=\color{black},
    % if not set, the frame-color may be changed on line-breaks
    % within not-black text (e.g. comments (green here))
    stepnumber=1,
    % the step between two line-numbers.
    % If it's 1, each line will be numbered
    stringstyle=\color{mymauve},     % string literal style
    tabsize=4,                       % sets default tabsize to 4 spaces
    % show the filename of files included with \lstinputlisting;
    % also try caption instead of title
    language = Python,
	showspaces = false,
	showtabs = false,
	showstringspaces = false,
	escapechar = ,
}

\def\ContinueLineNumber{\lstset{firstnumber=last}}
\def\StartLineAt#1{\lstset{firstnumber=#1}}
\let\numberLineAt\StartLineAt



\newcommand{\codeline}[1]{
	\alert{\texttt{#1}}
}


% Author and Course information
% This Document contains the information about this course.

% Authors of the slides
\author{Felix Döring, Felix Wittwer, Anton Obersteiner}

% Name of the Course
\institute{Python-Kurs}

% Fancy Logo 
\titlegraphic{\hfill\includegraphics[height=1.25cm]{../templates/fsr_logo_cropped}}



% Custom Bindings
% \newcommand{\codeline}[1]{
%	\alert{\texttt{#1}}
%}


% Presentation title
\title{Subprozesse in Python}
\date{\today}


\begin{document}

\maketitle

\begin{frame}{Gliederung}
        \begin{multicols}{2}
                \setbeamertemplate{section in toc}[sections numbered]
                \tableofcontents
        \end{multicols}
\end{frame}


\section{Grundlagen}
\begin{frame}[fragile]{Das Modul subprocess}
	Das Modul \alert{subprocess} erlaubt die Ausf\"uhrung externer Befehle und
	Skripte von einem Python Skript aus. Man kann sich die Funktionsweise \"ahnlich
	der eines Terminals vorstellen.
\end{frame}


\subsection{Eigenschaften}
\begin{frame}[fragile]{Eigenschaften}
	\begin{itemize}
		\item Subprozesse laufen \alert{asynchron}
		\item Sie laufen direkt auf dem System, nicht in einer Shell
		(wenn nicht anders festgelegt)
		\item Verf\"ugbare Programme h\"angen vom System ab, auf dem sie ausgef\"uhrt werden
		\item Der Aufruf ist allgemeing\"ultig, die Bibliothek verwandelt den
		Aufruf unter Windows in einen kompatiblen \texttt{CreateProcess()} String
	\end{itemize}
\end{frame}


\section{Konstanten}
\subsection{File Descriptoren}
\begin{frame}{File Descriptoren}
	\begin{description}
		\item[DEVNULL] Der systeminterne 'M\"ulleimer'
		\item[PIPE] Die Verbindung zwischen zwei Prozessen
		\item[STDOUT] Die Standardausgabe oder der laufende Prozess
	\end{description}
\end{frame}

\subsection{Exceptions}
\begin{frame}[fragile]{Exceptions}
	\begin{description}
		\item[\texttt{SubprocessError}] Der Standardfehler dieses Moduls
		\item[\texttt{TimeoutError}] Ein Timeout ist aufgetreten
		\item[\texttt{CalledProcessError}] Der Subprozess endete auf eine unerwartete Art
	\end{description}
\end{frame}


\section{Popen Klasse}
\begin{frame}[fragile]{Die Popen Klasse}
	Die \texttt{Popen} Klasse bildet die Basis des Moduls \texttt{subprocess}. \\
	Die Funktionssignatur sieht wie folgt aus:
	\lstinputlisting{resources/08_process/popen_call.py}
\end{frame}

\subsection{Wichtige Argumente}
\begin{frame}[fragile]{Einige wichtige Argumente}
	\begin{description}
		\item[\texttt{args}] Die aufzurufenden Argumente. Sollten vom Typ
			\alert{tuple} oder \alert{list} sein. Im Prinzip splittet man den
			Konsolenbefehl an den Leerzeichen: \\
			\texttt{ls -A *.md} entspricht \texttt{['ls', '-A', '*.md']}
		\item[\texttt{shell}] F\"uhrt den Befehl in einer Shell aus.
			Sollte \texttt{False} sein (Standard), sonst ist der Aufruf unsicher.
		\item[\texttt{stdout}] zusammen mit \texttt{stdin} und \texttt{stderr} die Input- und
			Output-Verbindungen des Subprozesses \\
			(hier sind \texttt{DEVNULL}, \texttt{PIPE} und \texttt{STDOUT} n\"utzlich)
		\item[\texttt{env}] Umgebungsvariablen des Kindprozesses. Standard ist ein
			Subset von \texttt{os.environ} (dem Python Prozess Environment)
		\item[\texttt{cwd}] Das Arbeitsverzeichnis des Subprozesses
	\end{description}
\end{frame}


\section{Popen Objekte}
\begin{frame}{Popen Objekte}
	Wenn Popen instanziiert wird, wird der darin enthaltene Prozess gestartet und
	das zur\"uckgegebene \texttt{Popen} Objekt enth\"alt Informationen \"uber den
	laufenden Prozess.
\end{frame}

\subsection{Informationen sammeln}
\begin{frame}[fragile]{Informationen sammeln}
	\begin{itemize}
		\item \texttt{process.args} \\
			Gibt die Argumente zur\"uck, mit denen der Prozess aufgerufen wurde.
		\item \texttt{obj.stdout}, \texttt{obj.stdin}, \texttt{obj.stderr} \\
			Input- und Output-Verbindungen, die beim Start gesetzt wurden
		\item \texttt{process.pid} \\
			Vom System zugewiesene \textit{Prozess ID}.
		\item \texttt{process.poll()} \\
			Pr\"uft, ob der Prozess beendet wurde. Gibt den \textit{R\"uckgabewert}
			des Prozesses zur\"uck oder \texttt{None}, wenn der Prozess noch l\"auft.
		\item \texttt{process.returncode} \\
			Der R\"uckgabewert von \texttt{process.poll()} (der R\"uckgabewert des Prozesses)
	\end{itemize}
\end{frame}

\subsection{Interaktion mit dem Prozess}
\begin{frame}[fragile]{Interaktion mit dem Prozess}
	\begin{itemize}
		\item \texttt{process.wait(timeout=None)} \\
			Wartet \textit{timeout} Sekunden auf die Terminierung des Prozesses
			(wartet unendlich lang, wenn timeout \textit{None} ist). \\
			Wirft nach Ablauf von \textit{timeout} eine \texttt{TimeoutExpired} Exception.
		\item \texttt{process.send\_signal(signal)} \\
			Sendet das Signal \textit{signal} an den Prozess (z.B. \texttt{SIGTERM}).
		\item \texttt{process.communicate(input=None, timeout=None)} \\
			Schreibt die Daten aus \textit{input} in den Standardinput (\texttt{stdin})
			des Prozesses (wenn stdin PIPE ist), wartet auf die Terminierung
			des Prozesses und liest die Daten, die der Prozess in \texttt{stdout}
			geschrieben hat (wenn stdout PIPE ist). \textit{timeout} funktioniert wie oben.
	\end{itemize}
\end{frame}

\begin{frame}[fragile]{Interaktion mit dem Prozess}
	\begin{itemize}
		\item \texttt{process.terminate()} \\
			Sendet ein Terminationssignal an den Prozess (\texttt{SIGTERM}).
		\item \texttt{process.kill()} \\
			Erzwingt die Beendigung des Prozesses (\texttt{SIGKILL}).
	\end{itemize}
\end{frame}

\subsection{Nutzung des Kontextmanagers}
\begin{frame}{Popen im Kontextmanager}
	Popen kann mit dem Kontextmanager verwendet werden (siehe \textit{File Handling}). \\
	Der Code daf\"ur w\"urde wie folgt aussehen:
	\lstinputlisting{resources/08_process/context_mgr.py}
\end{frame}

\section{N\"utzliche Funktionen}
\begin{frame}[fragile]{N\"utzliche Funktionen}
	\texttt{subprocess} enth\"alt einige Kurzfassungen f\"ur h\"aufig genutzte Arbeitsabl\"aufe.
	Intern rufen diese allerdings auch nur \texttt{Popen} auf.
\end{frame}

\begin{frame}[fragile]{run}
	Die wohl n\"utzlichste Funktion ist \texttt{subprocess.run()}:
	\lstinputlisting[lastline=2]{resources/08_process/defs.py}
	\begin{itemize}
		\item eingef\"uhrt in Python 3.5
		\item ruft durch \texttt{args} definierten Prozess auf
		\item wartet auf Beendigung des Prozesses (wenn \textit{timeout} \texttt{None} ist)
		\item gibt ein \texttt{CompletedProcess} Objekt zur\"uck
		\item \textbf{Beachte:} Einsatz eines unbenannten Aggregators
	\end{itemize}
\end{frame}

\begin{frame}[fragile]{CompletedProcess Objekt}
	Ist der Rückgabewert von \texttt{run()}, der zur\"uckgegeben wird, wenn der Prozess beendet wurde.
	Enth\"alt folgende Eigenschaften: \\[.25cm]
	\begin{description}
		\item[\texttt{args}] Argumente, mit denen der Prozess aufgerufen wurde.
		\item[\texttt{returncode}] R\"uckgabewert des Prozesses
		\item[\texttt{stdout}] Ausgabe des Prozesses
		\item[\texttt{stderr}] stderr Output des Prozesses \\[.5cm]
	\end{description}
	Au{\ss}erdem gibt es die Funktion \texttt{check\_returncode()}, die einen \texttt{CalledProcessError}
	beim Aufruf wirft, wenn der R\"uckgabewert nicht 0 ist.
\end{frame}

\begin{frame}[fragile]{call}
	\lstinputlisting[firstline=5, lastline=6]{resources/08_process/defs.py}
	\begin{itemize}
		\item Ruft einen Prozess auf, wartet auf Terminierung
		\item gibt Returncode des Prozesses zur\"uck \\[.5cm]
	\end{itemize}
	Diese Funktion macht das selbe, wie \\
	\hspace*{1.5cm}\texttt{run(...).returncode}, \\
	nur dass \textit{check} und \textit{input} nicht unterst\"utzt werden.
\end{frame}

\begin{frame}[fragile]{check\_call}
	Argumente entsprechen den Argumenten von \texttt{call()}.
	\begin{itemize}
		\item Ruft Prozess auf, wartet auf Terminierung
		\item gibt nichts zur\"uck, wenn Ausf\"uhrung erfolgreich (Returncode == 0)
		\item wenn Aufruf nicht erfolgreich, wird \texttt{CalledProcessError} geworfen \\[.5cm]
	\end{itemize}
	Diese Funktion macht das selbe, wie \\
	\hspace*{1.5cm}\texttt{run(..., check=True)}, \\
	nur dass \textit{input} nicht unterst\"utzt wird.
\end{frame}

\begin{frame}[fragile]{check\_output}
	\lstinputlisting[firstline=9]{resources/08_process/defs.py}
	\begin{itemize}
		\item f\"uhrt Kommando aus und gibt den Prozessoutput zur\"uck
		\item wirft ebenfalls \texttt{CalledProcessError}, wenn der Returncode nicht 0 ist \\[.5cm]
	\end{itemize}
	Diese Funktion macht das selbe, wie \\
	\hspace*{1.5cm}\texttt{run(..., check=True, stdout=PIPE).stdout}
\end{frame}

% nothing to do from here on
\end{document}
